\documentclass[11pt]{article}

\usepackage[backend=bibtex,natbib,style=numeric-comp,sorting=none,doi=false,isbn=false,url=false]{biblatex}
\addbibresource{../refs.bib}

\usepackage[T1]{fontenc}
\usepackage{lmodern}
\usepackage{amsmath,amssymb}
\usepackage{hyperref}
\usepackage[margin=1in]{geometry}

\newcommand{\Sym}{\mathrm{Sym}}
\newcommand{\mc}[1]{\mathcal{#1}}
\newcommand{\bR}{\mathbb{R}}
\newcommand{\Aeff}{\alpha'_{\mathrm{eff}}}

\title{Joining/Splitting in \texorpdfstring{$\Sym^M(\bR^8)$}{SymM(R8)}:\\
Large-\texorpdfstring{$M$}{M} Counting with Macroscopic Cycles,\\
and a Sphere-Target First-Order Worldsheet Analysis}
\author{}
\date{}

\begin{document}
\maketitle

\section{Aim and scope}
This note presents a self-contained comparison between symmetric-product-orbifold three-point data and the curved first-order worldsheet description of joining/splitting processes. The main goals are:
\begin{itemize}
\item State the exact finite-$M$ orbifold three-point formula and separate two inequivalent large-$M$ regimes.
\item Work out the macroscopic-cycle limit $w_i\sim M$ with $w_1+w_2+w_3=M$, including explicit combinatorial asymptotics and the reduced object used for controlled $1/M$ organization.
\item Derive sphere three-point localization directly from the curved first-order Gomis--Ooguri path integral, showing in detail why branched covers appear.
\item Incorporate the global BRST analysis of the curved target and explain how it changes the sphere three-point comparison and the interpretation of the $1/M$ story.
\end{itemize}

The notation is:
\begin{equation}
\mc C_M:=\Sym^M(\bR^8),\qquad M\in \mathbb{Z}_{>0},
\end{equation}
and $w_i$ are cycle lengths of the three external twist operators.

\section{NRCS cylinder scaling, large-\texorpdfstring{$N$}{N}, and the IR orbifold regime}
\label{sec:nrcs-cylinder-ir}
Consider NRCS with longitudinal target coordinates $(x^0,x^1)$, where $x^0$ is noncompact time and
\begin{equation}
x^1\sim x^1+2\pi R.
\end{equation}
The cylinder radius is $R$, and we define the dimensionless radius
\begin{equation}
\rho:=\frac{R}{\sqrt{\Aeff}}.
\end{equation}

In this section $N$ denotes the rank of the dual $U(N)$ gauge theory. We use the NCOS/NRCS convention
\begin{equation}\label{eq:alphaeff-gym-n}
\Aeff=\frac{N^2}{g_{\mathrm{YM}}^2},
\end{equation}
so $g_{\mathrm{YM}}$ has mass dimension one in $1+1$ dimensions. (Some references use a $2\pi$-shifted normalization; this only changes fixed numerical factors and not the large-$N$ scaling logic \cite{klebanov200011dimensionalncosun}.)

Define
\begin{equation}\label{eq:kappa-def}
\kappa:=g_{\mathrm{YM}}R.
\end{equation}
Using \eqref{eq:alphaeff-gym-n},
\begin{equation}\label{eq:kappa-rho}
\kappa=N\rho=N\frac{R}{\sqrt{\Aeff}}.
\end{equation}
So there are two inequivalent large-$N$ scalings:
\begin{equation}\label{eq:largeN-two-scalings}
\text{(A) fixed }\rho:\ \kappa=N\rho\to\infty,
\qquad
\text{(B) fixed }\kappa:\ \rho=\kappa/N\to 0.
\end{equation}
In this section we use (A): fix $R/\sqrt{\Aeff}$ and scale only $N$.
Then
\begin{equation}\label{eq:oneoverN-oneoverkappa}
\frac{1}{N}=\frac{\rho}{\kappa},
\qquad
\frac{1}{\kappa}=\frac{1}{g_{\mathrm{YM}}R}.
\end{equation}
So at fixed $\rho$, a $1/N$ expansion is exactly a $1/(g_{\mathrm{YM}}R)$ expansion up to the fixed factor $\rho$.

By ``finite-$N$ NRCS'' we mean a hypothetical exact (nonperturbative) definition that captures the full finite-$N$ gauge theory. That exact formulation is not currently available. By contrast, perturbative NRCS is organized as a large-$N$ expansion (equivalently in the NCOS language, an expansion around weak effective string coupling), so it automatically probes the large-$N$ regime \cite{gomis2000nonrelativisticclosedstringtheory,klebanov200011dimensionalncosun}.

Equation \eqref{eq:alphaeff-gym-n} gives
\begin{equation}
\frac{1}{\sqrt{\Aeff}}=\frac{g_{\mathrm{YM}}}{N}.
\end{equation}
For any excitation with finite dimensionless cylinder energy $\varepsilon:=ER$,
\begin{equation}
\frac{E}{g_{\mathrm{YM}}}
=
\frac{ER}{g_{\mathrm{YM}}R}
=
\frac{\varepsilon}{\kappa}.
\end{equation}
Therefore in scaling (A), perturbative states with $\varepsilon=O(1)$ satisfy
\begin{equation}
\frac{E}{g_{\mathrm{YM}}}=O(N^{-1})\ll 1,
\end{equation}
which is an IR scale of the gauge theory. This is the regime where the matrix-string description flows to the symmetric-product orbifold CFT \cite{dijkgraaf1997matrixstring}.

To avoid notation clash with the rest of the note: below we continue to use $M$ for orbifold copy number in $\Sym^M(\bR^8)$, and the comparison with this section uses $M=N$.

Once the IR description is a CFT on a cylinder of radius $R$, changing $R$ rescales the Hamiltonian but not local OPE data:
\begin{equation}
H(R)=\frac{1}{R}\left(L_0+\bar L_0-\frac{c}{12}\right),
\qquad
P(R)=\frac{1}{R}\left(L_0-\bar L_0\right).
\end{equation}
So local CFT data (for example OPE coefficients and normalized local correlators) are independent of the overall scale $R$, while energies in physical units scale as $1/R$.

The same pattern appears directly on the NRCS side. Importing the continuum lightcone NRCS spectrum (rather than deriving it in this note) \cite{gomis2000nonrelativisticclosedstringtheory}, for winding $w>0$ one has
\begin{equation}\label{eq:nrcs-spectrum-cylinder}
p_0=
\frac{wR}{2\Aeff}
 +\frac{\Aeff\,k_\perp^2}{2wR}
 +\frac{N_L+N_R}{wR},
\qquad
N_L-N_R=wn.
\end{equation}
Here $k_\perp$ is transverse momentum, $N_L,N_R$ are left/right oscillator levels, and $n$ is the integer momentum mode along $x^1$.
Subtracting the nonrelativistic rest energy,
\begin{equation}
\Delta E:=p_0-\frac{wR}{2\Aeff}
=
\frac{1}{R}
\left[
\frac{(k_\perp\sqrt{\Aeff})^2}{2w}
 +\frac{N_L+N_R}{w}
\right].
\end{equation}
Thus $R$ appears only as an overall $1/R$ factor in excitation energies.
Equivalently,
\begin{equation}\label{eq:deltaE-over-gYM}
\frac{\Delta E}{g_{\mathrm{YM}}}
=
\frac{1}{\kappa}
\left[
\frac{(k_\perp\sqrt{\Aeff})^2}{2w}
 +\frac{N_L+N_R}{w}
\right].
\end{equation}

In the macroscopic winding sector $w=xN$ with fixed $0<x\le 1$ and fixed $\kappa=g_{\mathrm{YM}}R$, \eqref{eq:kappa-def} gives
\begin{equation}
\Delta E
=
\frac{g_{\mathrm{YM}}}{\kappa N}
\left[
\frac{(k_\perp\sqrt{\Aeff})^2}{2x}
 +\frac{N_L+N_R}{x}
\right].
\end{equation}
Using $\kappa=N\rho$, this is equivalently
\begin{equation}\label{eq:deltaE-macro-rho}
\Delta E
=
\frac{g_{\mathrm{YM}}}{N^2\rho}
\left[
\frac{(k_\perp\sqrt{\Aeff})^2}{2x}
+\frac{N_L+N_R}{x}
\right].
\end{equation}
Hence in fixed-$\rho$ large-$N$:
\begin{itemize}
\item generic winding sectors with $w=O(1)$ have $\Delta E/g_{\mathrm{YM}}=O(1/N)$ from \eqref{eq:deltaE-over-gYM},
\item macroscopic sectors with $w\sim N$ have $\Delta E/g_{\mathrm{YM}}=O(1/N^2)$ from \eqref{eq:deltaE-macro-rho}.
\end{itemize}
Both are IR in gauge-theory units, consistent with landing on the symmetric-product orbifold regime. The overall radius $R$ enters only through the dimensionless ratio $\rho=R/\sqrt{\Aeff}$.

This also matches the DVV organization of matrix string theory. The deformation
\begin{equation}
S_{\mathrm{int}}=\lambda_{\mathrm{DVV}}\int d^2z\,\mc O_{(2)}
\end{equation}
has $h(\mc O_{(2)})=\bar h(\mc O_{(2)})=3/2$, so $\lambda_{\mathrm{DVV}}$ has length dimension. On a cylinder of radius $R$, the dimensionless coupling is
\begin{equation}\label{eq:dvv-dimensionless}
\widehat\lambda_{\mathrm{DVV}}:=\frac{\lambda_{\mathrm{DVV}}}{R}.
\end{equation}
In the standard MST dictionary one has
\begin{equation}
\lambda_{\mathrm{DVV}}=\frac{c_{\mathrm{DVV}}}{g_{\mathrm{YM}}},
\end{equation}
with convention-dependent $c_{\mathrm{DVV}}=O(1)$ determined by operator normalization. Therefore
\begin{equation}\label{eq:dvv-kappa-relation}
\widehat\lambda_{\mathrm{DVV}}=\frac{c_{\mathrm{DVV}}}{g_{\mathrm{YM}}R}
=\frac{c_{\mathrm{DVV}}}{\kappa},
\qquad
\frac{1}{N}=\frac{\rho}{c_{\mathrm{DVV}}}\,\widehat\lambda_{\mathrm{DVV}}.
\end{equation}
So in perturbative NRCS scaling (fixed $\rho$), the NRCS $1/N$ expansion and the DVV/MST expansion are the same parametric expansion.

\section{Orbifold three-point data at finite \texorpdfstring{$M$}{M}}
At the orbifold point, define the normalized single-cycle class operator
\begin{equation}\label{eq:class-op}
\widehat{\sigma}_{[w]}
:=
\frac{1}{\sqrt{|[w]|_M}}
\sum_{g\in [w]}\sigma_g,
\qquad
|[w]|_M=\frac{M!}{w(M-w)!}.
\end{equation}
Here $[w]\subset S_M$ is the conjugacy class of $w$-cycles.

For the connected three-point contribution associated with degree-$d$ covers of genus $g$,
\begin{equation}\label{eq:C123-def}
C_{123}^{(g)}(M)
:=
\left\langle
\widehat{\sigma}_{[w_1]}(0)\,
\widehat{\sigma}_{[w_2]}(1)\,
\widehat{\sigma}_{[w_3]}(\infty)
\right\rangle_{g}^{\mathrm{conn}},
\end{equation}
the standard covering-space decomposition gives
\begin{equation}\label{eq:C123-exact}
C_{123}^{(g)}(M)
=
\frac{M!}{(M-d)!}\,
\frac{
H_d^{(g)}(w_1,w_2,w_3)\,
\mc C_{\mathrm{seed}}^{(g)}(w_1,w_2,w_3)
}{
\sqrt{|[w_1]|_M\,|[w_2]|_M\,|[w_3]|_M}
}.
\end{equation}
Definitions:
\begin{itemize}
\item $d$ is the number of active copies (equivalently covering degree).
\item $H_d^{(g)}$ is the Hurwitz count $\sum_f 1/|\mathrm{Aut}(f)|$ over covers with the prescribed ramification.
\item $\mc C_{\mathrm{seed}}^{(g)}$ is the local seed-theory correlator on the cover.
\end{itemize}

For three branch points, Riemann--Hurwitz gives
\begin{equation}\label{eq:RH}
2-2g
=
2d-\sum_{i=1}^3(w_i-1),
\qquad
d=\frac{w_1+w_2+w_3-1}{2}-g.
\end{equation}

\section{Two inequivalent large-\texorpdfstring{$M$}{M} limits}
\subsection{Dilute-cycle regime: fixed \texorpdfstring{$w_i$}{wi} as \texorpdfstring{$M\to\infty$}{M->infty}}
If $w_i$ and $d$ are fixed while $M\to\infty$, then
\begin{equation}
\frac{M!}{(M-d)!}
=
M^d\left[
1-\frac{d(d-1)}{2M}+O(M^{-2})
\right],
\end{equation}
\begin{equation}
|[w]|_M
=
\frac{M^w}{w}
\left[
1-\frac{w(w-1)}{2M}+O(M^{-2})
\right].
\end{equation}
Substituting into \eqref{eq:C123-exact} yields
\begin{equation}\label{eq:dilute-expansion}
\begin{aligned}
C_{123}^{(g)}(M)
=&
M^{-g-\frac12}\,
\sqrt{w_1w_2w_3}\,
H_d^{(g)}\,
\mc C_{\mathrm{seed}}^{(g)}
\\
&\times
\left[
1+\frac{1}{M}
\left(
\frac14\sum_{i=1}^3 w_i(w_i-1)-\frac12 d(d-1)
\right)
+O(M^{-2})
\right].
\end{aligned}
\end{equation}
This is the usual power-series $1/M$ expansion.

\subsection{Macroscopic-cycle regime: \texorpdfstring{$w_i\sim M$}{wi~M} and \texorpdfstring{$w_1+w_2+w_3=M$}{sum wi=M}}
Now impose the scaling that you asked for:
\begin{equation}\label{eq:macro-scaling}
w_i=x_i M+O(1),
\qquad
0<x_i<1,
\qquad
x_1+x_2+x_3=1.
\end{equation}
For genus $g=0$, \eqref{eq:RH} gives
\begin{equation}\label{eq:d-macro}
d=\frac{M-1}{2}.
\end{equation}
Define the pure combinatorial prefactor
\begin{equation}\label{eq:Acomb-def}
\mc A_{\mathrm{comb}}(M;\{w_i\})
:=
\frac{M!}{(M-d)!}\,
\frac{1}{\sqrt{|[w_1]|_M|[w_2]|_M|[w_3]|_M}},
\end{equation}
so
\begin{equation}\label{eq:C-macro-factorized}
C_{123}^{(0)}(M)
=
\mc A_{\mathrm{comb}}(M;\{w_i\})\,
H_{(M-1)/2}^{(0)}(w_1,w_2,w_3)\,
\mc C_{\mathrm{seed}}^{(0)}(w_1,w_2,w_3).
\end{equation}

Using Stirling in this scaling regime gives
\begin{equation}\label{eq:Acomb-stirling}
\log \mc A_{\mathrm{comb}}(M;\{x_i\})
=
\frac{M}{2}\,\Psi(x_1,x_2,x_3)
+\frac{3}{2}\log M
+O(1),
\end{equation}
with
\begin{equation}\label{eq:Psi-def}
\Psi(x_1,x_2,x_3)
:=
\log 2+\sum_{i=1}^3 (1-x_i)\log(1-x_i).
\end{equation}
Therefore
\begin{equation}\label{eq:Acomb-asymptotic}
\mc A_{\mathrm{comb}}(M;\{x_i\})
=
\mc N(x_1,x_2,x_3)\,
M^{3/2}\,
\exp\!\left[\frac{M}{2}\Psi(x_1,x_2,x_3)\right]
\left[1+O(M^{-1})\right],
\end{equation}
where $\mc N(x_1,x_2,x_3)$ is an $M$-independent prefactor from subleading Stirling terms.

This is the key point: in the macroscopic regime, the full correlator is \emph{not} a naive power series in $1/M$. The cycle fractions $x_i$ enter already at exponential order through \eqref{eq:Psi-def}, and also through large-degree Hurwitz data and the large-$w_i$ behavior of $\mc C_{\mathrm{seed}}^{(0)}$.

A useful reduced object is
\begin{equation}\label{eq:reduced-amplitude}
\widetilde C_{123}^{(0)}(M)
:=
e^{-\frac{M}{2}\Psi(x_1,x_2,x_3)}\,M^{-3/2}\,C_{123}^{(0)}(M),
\end{equation}
for which one may then attempt a genuine $1/M$ expansion if
\begin{equation}
H_{(M-1)/2}^{(0)}(w_1,w_2,w_3)\,
\mc C_{\mathrm{seed}}^{(0)}(w_1,w_2,w_3)
\end{equation}
admits such an expansion in the same scaling limit.

\section{DVV deformation versus large-\texorpdfstring{$M$}{M} counting}
The matrix-string deformation is
\begin{equation}\label{eq:dvv}
S=S_{\mathrm{orb}}+\lambda_{\mathrm{DVV}}\int d^2z\,\mc O_{(2)}(z,\bar z),
\qquad
h(\mc O_{(2)})=\bar h(\mc O_{(2)})=\frac{3}{2},
\end{equation}
and first order gives
\begin{equation}\label{eq:dvv-first}
\delta C_{123}
=
-\lambda_{\mathrm{DVV}}
\int d^2z\,
\left\langle
\widehat{\sigma}_{[w_1]}\widehat{\sigma}_{[w_2]}\widehat{\sigma}_{[w_3]}
\mc O_{(2)}(z,\bar z)
\right\rangle_{\mathrm{orb}}.
\end{equation}

The power of $M$ in a normalized correlator is controlled by the combinatorics in \eqref{eq:C123-exact}, not by \eqref{eq:dvv} alone. In general:
\begin{itemize}
\item ``order in $\lambda_{\mathrm{DVV}}$'' and ``order in $1/M$'' are logically distinct,
\item and in the macroscopic regime \eqref{eq:macro-scaling}, even the undeformed correlator is not a pure $1/M$ series before extracting the entropic factor.
\end{itemize}
To connect this with perturbative NRCS scaling, use the dimensionless DVV coupling
\begin{equation}
\widehat\lambda_{\mathrm{DVV}}:=\frac{\lambda_{\mathrm{DVV}}}{R},
\end{equation}
and the relation \eqref{eq:dvv-kappa-relation}
\begin{equation}
\widehat\lambda_{\mathrm{DVV}}=\frac{c_{\mathrm{DVV}}}{\kappa},
\qquad
c_{\mathrm{DVV}}=O(1)\ \text{(convention-dependent)}.
\end{equation}
Then in the perturbative NRCS scaling of Section \ref{sec:nrcs-cylinder-ir} with $M=N$ and fixed $\rho=R/\sqrt{\Aeff}$,
\begin{equation}
\frac{1}{M}=\frac{1}{N}=\frac{\rho}{\kappa}
=
\frac{\rho}{c_{\mathrm{DVV}}}\,\widehat\lambda_{\mathrm{DVV}}.
\end{equation}
So the NRCS $1/N$ expansion and the DVV/MST expansion are the same \emph{parametric} expansion in $1/\kappa=1/(g_{\mathrm{YM}}R)$.

The practical distinction is controlled by $\rho$:
\begin{itemize}
\item if $\rho=O(1)$, $1/M$ and $\widehat\lambda_{\mathrm{DVV}}$ are comparable at fixed $\kappa$,
\item if $\rho\ll 1$, $1/M$ effects are suppressed relative to DVV insertions,
\item if $\rho\gg 1$, $1/M$ effects are enhanced relative to DVV insertions.
\end{itemize}
Therefore a safe organization is a double expansion
\begin{equation}
C_{123}
=
\sum_{n,m\ge 0}
a_{n,m}\!\left(\rho;\,x_i,\ldots\right)\,
\widehat\lambda_{\mathrm{DVV}}^{\,n}\,M^{-m},
\end{equation}
with explicit $\rho$-dependence kept in the coefficients. Only after fixing $\rho$ does this collapse to a single $1/\kappa$ expansion.

\section{Worldsheet three-point from the curved first-order action on \texorpdfstring{$S^2$}{S2}}
\subsection{Action and localization equations}
Use the same curved first-order matter system as in \texttt{CURVED\_FIRST\_ORDER\_STRING.tex}, specialized to
\begin{equation}
ds^2_{S^2}=e^{2\omega(u,\bar u)}\,du\,d\bar u,
\qquad
e^{2\omega(u,\bar u)}=\frac{4}{(1+|u|^2)^2},
\end{equation}
with superderivatives
\begin{equation}
D=\partial_\vartheta+\vartheta\partial,\qquad
\bar D=\partial_{\bar\vartheta}+\bar\vartheta\bar\partial,
\end{equation}
and action
\begin{equation}\label{eq:worldsheet-action}
\begin{aligned}
S_{\mathrm{ws}}
=
\frac{1}{2\pi}\int d^2z\,d\vartheta\,d\bar\vartheta
\Big[
\mathbf B\,\bar D\mathbf X^+
+\bar{\mathbf B}\,D\mathbf X^-
+\frac{e^{2\omega(\mathbf X^+,\mathbf X^-)}}{4\Aeff}\,
D\mathbf X^+\,\bar D\mathbf X^-
\Big]
+S_\perp.
\end{aligned}
\end{equation}

Consider the three-point function
\begin{equation}\label{eq:A3-pathintegral}
\mc A_3
=
\int [D\Phi]\,
e^{-S_{\mathrm{ws}}[\Phi]}\,
\prod_{i=1}^3 \mc V_{w_i}(z_i,\bar z_i),
\end{equation}
where $\Phi$ denotes all worldsheet fields and
\begin{equation}
\mc V_{w_i}=\sigma_{w_i}^{\mathrm{long}}\cdot \mc V_i^\perp
\end{equation}
is a longitudinal twist insertion times a transverse operator.

\paragraph{Step 1: integrating out first-order multiplets.}
The longitudinal part of \eqref{eq:worldsheet-action} is first order:
\begin{equation}
S_{\mathrm{long}}
=
\frac{1}{2\pi}\int_\Sigma
\Big(
\mathbf B\,\bar D\mathbf X^+
\bar{\mathbf B}\,D\mathbf X^-
\cdots
\Big).
\end{equation}
The omitted terms are fermionic/background couplings from \eqref{eq:worldsheet-action} that do not contain $\mathbf B$ or $\bar{\mathbf B}$, so they do not modify the $\delta$-functional constraints derived below.
At fixed $\mathbf X^\pm$, the $\mathbf B$ integral factorizes mode-by-mode. For a mode coefficient $b_n$ one gets
\begin{equation}
\int db_n\,e^{-b_n\,(\bar D\mathbf X^+)_n}
\propto
\delta\!\big((\bar D\mathbf X^+)_n\big),
\end{equation}
so the full functional integral gives
\begin{equation}\label{eq:delta-functional}
\int [D\mathbf B]\,
\exp\!\left[
-\frac{1}{2\pi}\int \mathbf B\,\bar D\mathbf X^+
\right]
\propto
\delta[\bar D\mathbf X^+]\,
\big(\det{}'\bar D\big)^{-1},
\end{equation}
and similarly for $\bar{\mathbf B}$. The prime on $\det{}'$ reminds us that zero-mode conventions are separated from nonzero-mode determinants.

Hence the support of the path integral is
\begin{equation}\label{eq:chiral-constraints}
\bar D\mathbf X^+=0,\qquad D\mathbf X^-=0.
\end{equation}
On the Euclidean contour $\mathbf X^-=\overline{\mathbf X^+}$ this gives
\begin{equation}
X^+=f(z),\qquad X^-=\overline{f(\bar z)},
\end{equation}
with $f$ holomorphic on $\Sigma$ away from operator insertions.

\paragraph{Step 2: monodromy from twist insertions.}
The longitudinal twist field of length $w_i$ cyclically permutes $w_i$ copies, so going once around $z_i$ moves one to the next sheet and returns only after $w_i$ turns. In local uniformizing coordinate $t_i$,
\begin{equation}
z-z_i=t_i^{\,w_i},
\end{equation}
and $f$ is single-valued in $t_i$. Equivalently, near $z_i$,
\begin{equation}\label{eq:branch-local}
f(z)=u_i+a_i(z-z_i)^{w_i}+O((z-z_i)^{w_i+1}),
\end{equation}
so
\begin{equation}
f'(z)\sim (z-z_i)^{w_i-1},
\end{equation}
and the ramification index is $r_i=w_i-1$.

\paragraph{Step 3: global classification.}
Global solutions of \eqref{eq:chiral-constraints} compatible with \eqref{eq:branch-local} are exactly branched holomorphic covers
\begin{equation}
f:\Sigma\to S^2
\end{equation}
with branch indices $r_i=w_i-1$ and degree $d$ constrained by \eqref{eq:RH}. For $\Sigma=S^2$, every holomorphic map to $S^2$ is rational, and once $z_i$ are fixed to $(0,1,\infty)$ there is no residual continuous modulus for three marked points. Therefore localization means: the functional integral over longitudinal fields collapses to a discrete sum over the corresponding Hurwitz class, with fluctuation determinants and Jacobians supplying the weight of each cover.

\subsection{Localized form of the amplitude}
The localized expression from \eqref{eq:A3-pathintegral} is
\begin{equation}\label{eq:A3-localized}
\mc A_3
=
\sum_{f\in \mathrm{Hurwitz}_0(d;w_1,w_2,w_3)}
\frac{1}{|\mathrm{Aut}(f)|}\,
\exp[-S_{\mathrm{cl}}[f]]\,
\mc Z^{\mathrm{long}}_{1\text{-}\mathrm{loop}}[f]\,
\mc C_\perp[f].
\end{equation}
Here:
\begin{itemize}
\item $S_{\mathrm{cl}}[f]$ is the classical action of \eqref{eq:worldsheet-action} evaluated on $f$.
\item $\mc Z^{\mathrm{long}}_{1\text{-}\mathrm{loop}}[f]$ is the renormalized determinant/Jacobian from longitudinal fluctuations and ghosts.
\item $\mc C_\perp[f]$ is the transverse correlator on the same cover.
\end{itemize}

Equation \eqref{eq:A3-localized} is derived directly from the first-order action and does not assume the final Hurwitz-only form.

\section{Sphere target after the global BRST analysis}
The note \texttt{BRST-curved.tex} shows that the relevant distinction is between:
\begin{itemize}
\item the \emph{bosonic} curved $(\beta,\gamma)$ subsystem, where the usual obstruction is controlled by $c_1(TC)$ \cite{nekrasov2005curvedbetagammasystems}, and
\item the \emph{full super} longitudinal multiplet $(\gamma,\beta,\psi^+,\rho)$, where that obstruction is cancelled.
\end{itemize}
For a target curve $C$, the bosonic topological class is
\begin{equation}\label{eq:c1-s2}
\int_C c_1(TC)=\chi(C)=2-2g_C,
\end{equation}
and for $C=S^2$ this gives $\int_{S^2}c_1(TS^2)=2$.

The local patching formulas summarized in \texttt{BRST-curved.tex} are
\begin{equation}\label{eq:super-patching-summary}
\gamma_b=f_{ba}(\gamma_a),\quad
\psi_b^+=f'_{ba}(\gamma_a)\psi_a^+,\quad
\rho_b=\bigl(f'_{ba}(\gamma_a)\bigr)^{-1}\rho_a,\quad
\beta_b=\bigl(f'_{ba}(\gamma_a)\bigr)^{-1}\beta_a-\frac{f''_{ba}(\gamma_a)}{\bigl(f'_{ba}(\gamma_a)\bigr)^2}\psi_a^+\rho_a.
\end{equation}
With $\Gamma_a:=2\partial_{u_a}\omega_a$ and
\begin{equation}
\Pi_a:=\beta_a-\Gamma_a\rho_a\psi_a^+,
\end{equation}
one gets tensorial overlap behavior
\begin{equation}\label{eq:covariant-momentum-summary}
\Pi_b=\bigl(f'_{ba}\bigr)^{-1}\Pi_a.
\end{equation}
Now define $J_a:=:\rho_a\psi_a^+:$. In one target complex dimension, $J_a^2=0$, so
\begin{equation}\label{eq:supercurrent-global-summary}
:\psi_a^+\Pi_a: = :\psi_a^+\beta_a:,
\end{equation}
and the stress tensor reduces to
\begin{equation}\label{eq:stress-global-summary}
T_{+-}^{(a)}
=
-:\beta_a\partial\gamma_a:
-\frac12:\rho_a\partial\psi_a^+:
+\frac12:(\partial\rho_a)\psi_a^+:.
\end{equation}
These are the flat-form local operators, now shown to patch globally on a curve. Quantum mechanically, \texttt{BRST-curved.tex} further argues that the bosonic determinant line is cancelled by the fermionic partner determinant line, consistent with the chiral-de-Rham perspective \cite{msv1999chiralderham,benzvi2008supersymmetrycdr}. Therefore, for the \emph{full} super longitudinal multiplet, there is no residual curved-$(\beta,\gamma)$ anomaly factor to insert in the amplitude.

For the dilaton-like term,
\begin{equation}\label{eq:dilaton-term}
S_\Phi[f]
=
\frac{1}{4\pi}\int_\Sigma \sqrt{g_{\mathrm{ws}}}\,R^{(2)}_{\mathrm{ws}}\,\Phi(f).
\end{equation}
Decompose
\begin{equation}
\Phi(f)=\Phi_0+\delta\Phi(f),
\end{equation}
then for $\Sigma=S^2$ (so $\int_\Sigma \sqrt{g_{\mathrm{ws}}}R^{(2)}_{\mathrm{ws}}=8\pi$),
\begin{equation}\label{eq:Sphi-sphere}
S_\Phi[f]
=
2\Phi_0
+\frac{1}{4\pi}\int_\Sigma \sqrt{g_{\mathrm{ws}}}\,R^{(2)}_{\mathrm{ws}}\,\delta\Phi(f).
\end{equation}
If $\Phi$ is constant on the target, the second term vanishes and $S_\Phi[f]=2\Phi_0$ exactly.
If $\Phi$ is nonconstant, that second term is map-dependent through $f$ and therefore changes relative weights among covers.

Hence the saddle weight used in the localized sum is
\begin{equation}\label{eq:anomaly-corrected-weight}
\mc W_f
=
\exp\!\Big[-S_{\mathrm{cl}}[f]-S_\Phi[f]\Big]\,
\mc Z^{\mathrm{long}}_{1\text{-}\mathrm{loop}}[f],
\end{equation}
and \eqref{eq:A3-localized} is
\begin{equation}\label{eq:A3-corrected}
\mc A_3
=
\sum_{f}\frac{1}{|\mathrm{Aut}(f)|}\,\mc W_f\,\mc C_\perp[f].
\end{equation}
So after incorporating the global BRST result, the unresolved input is no longer a topological anomaly coefficient. The unresolved input is the dynamical map dependence in $S_{\mathrm{cl}}[f]$, possible map dependence in $S_\Phi[f]$, and the renormalized determinant $\mc Z^{\mathrm{long}}_{1\text{-}\mathrm{loop}}[f]$.
The same BRST analysis also implies that the finite-$\Aeff$ interaction changes nonchiral dynamics but does not change the chiral BRST complex. In this localized three-point problem, that means the remaining uncertainty enters through dynamical weights in \eqref{eq:anomaly-corrected-weight}, not through a failure of global BRST definition.

\section{Comparison with orbifold side}
The orbifold formula \eqref{eq:C123-exact} is exact at the free point.
The worldsheet formula derived from the action is \eqref{eq:A3-corrected}.
Relative to the previous version, one important term is removed: there is no extra factor $e^{-2\kappa_{\mathrm{anom}}d}$ (with $\kappa_{\mathrm{anom}}$ the would-be bosonic anomaly coefficient) from a residual sphere anomaly of the longitudinal supermultiplet.

The two descriptions match structurally if the following conditions hold:
\begin{itemize}
\item The map-dependent part in \eqref{eq:anomaly-corrected-weight} is absent or can be re-expressed in orbifold data. Concretely, this means the factor from $S_{\mathrm{cl}}[f]$, $S_\Phi[f]$, and $\mc Z^{\mathrm{long}}_{1\text{-}\mathrm{loop}}[f]$ should be independent of $f$ inside a fixed Hurwitz class.
\item The transverse factor $\mc C_\perp[f]$ matches the seed correlator data entering $\mc C_{\mathrm{seed}}^{(g)}$.
\item External-state normalization is matched between BRST vertex operators and class-normalized orbifold operators \eqref{eq:class-op}.
\end{itemize}

Under those conditions one recovers the expected Hurwitz organization. Without those conditions, the $S^2$ three-point comparison remains an honest \emph{conditional} statement:
\begin{equation}
\mc A_3
\propto
\sum_f\frac{1}{|\mathrm{Aut}(f)|}\,
\big[\text{map-dependent correction}\big]\,
\mc C_\perp[f].
\end{equation}
So coefficient-level equality with the orbifold answer is not yet proved unless the bracketed correction is controlled.

Equivalently:
\begin{itemize}
\item If $S_{\mathrm{cl}}[f]$, $S_\Phi[f]$, and $\mc Z^{\mathrm{long}}_{1\text{-}\mathrm{loop}}[f]$ are cover-independent within fixed branching data, then the worldsheet result differs from the orbifold one only by an overall normalization that can be absorbed into coupling/operator conventions.
\item If any of those factors is genuinely cover-dependent, worldsheet weights are no longer pure Hurwitz weights, and exact coefficient-by-coefficient equality with \eqref{eq:C123-exact} fails unless one adds the corresponding correction on the orbifold side.
\end{itemize}

\section{What is currently solid for macroscopic scaling}
For $w_i\sim M$ with $w_1+w_2+w_3=M$:
\begin{itemize}
\item The exact orbifold expression is \eqref{eq:C-macro-factorized}.
\item The combinatorial prefactor has explicit asymptotics \eqref{eq:Acomb-stirling}--\eqref{eq:Acomb-asymptotic}, so the $x_i=w_i/M$ fractions enter nontrivially already at leading order.
\item Therefore the full correlator is not a naive power series in $1/M$ before extracting the entropic factor \eqref{eq:reduced-amplitude}.
\item On the worldsheet side, the first-order action gives localization to branched covers \eqref{eq:A3-corrected}, and there is no additional unknown anomaly factor of the form $e^{-\kappa M}$ from the longitudinal supermultiplet.
\item The remaining uncertainty for coefficient-level matching is dynamical map dependence in $S_{\mathrm{cl}}[f]$, $S_\Phi[f]$, and $\mc Z^{\mathrm{long}}_{1\text{-}\mathrm{loop}}[f]$, not a residual curved-$(\beta,\gamma)$ topological obstruction.
\end{itemize}

\section{Do we need grand-canonical variables here?}
For the torus free-energy problem, grand-canonical variables are essential because one sums over all cycle multiplicities (including disconnected sectors) and then extracts fixed-$M$ data via plethystic operations.

For the present sphere three-point observable, we insert three \emph{fixed} operators of prescribed cycle lengths $(w_1,w_2,w_3)$ and fixed rank $M$. This is already a canonical correlator:
\begin{equation}
\left\langle
\widehat{\sigma}_{[w_1]}\widehat{\sigma}_{[w_2]}\widehat{\sigma}_{[w_3]}
\right\rangle_{\Sym^M}^{\mathrm{conn}},
\end{equation}
so no grand-canonical fugacity is required to define or compute that single quantity.

Grand-canonical generating functions can still be useful if one wants to organize an infinite family of three-point coefficients over varying $(w_1,w_2,w_3)$, but that is optional bookkeeping, not a structural necessity for a single fixed 3-point function.

\section{Key references used}
\begin{itemize}
\item Matrix-string interaction logic and contact terms: \cite{dijkgraaf1997matrixstring,dijkgraaf2003matrixstringcontactterms}.
\item 1+1 NCOS/SYM parameter dictionary and large-$N$ scaling: \cite{klebanov200011dimensionalncosun}.
\item Symmetric-orbifold covering-map correlators and large-$M$ combinatorics: \cite{arutyunov1998orbifoldscattering,lunin2001correlationorbifolds,lunin2002threepointorbifolds,pakman2009diagramssymmetricorbifold,pakman2009extremalhurwitz}.
\item First-order NCOS/GO worldsheet system: \cite{gomis2000nonrelativisticclosedstringtheory}.
\item Curved-$(\beta,\gamma)$ anomaly framework and chiral-de-Rham cancellation structure: \cite{nekrasov2005curvedbetagammasystems,msv1999chiralderham,benzvi2008supersymmetrycdr}.
\item Global BRST and curved-target patching formulas used in this note are implemented in \texttt{BRST-curved.tex}.
\end{itemize}

\printbibliography

\end{document}
